


\section{Two creditors, policy consistency and financing refusal}
In a finite one-period model, policy consistency can be checked by enumeration.
Specify a candidate policy $\pi$, the distribution of the terminal estate under
that policy, the priority of each claim, and a pricing measure and discount
factor. Allocate the estate once, calculate each lender's expected discounted
receipt, and use those receipts as candidate issue proceeds. At those fixed
quoted proceeds, compare the issuer's payoffs from every feasible policy.
Retain the candidate only when its conjectured policy is a best response and
all funding and participation constraints hold. Ties must be retained. If no
pure-policy pair passes, report that result for the declared class; neither a
balancing loan nor an assumed mixed equilibrium repairs it automatically.

For a concrete local example, existing senior debt has face 40 and newly issued
junior debt has face 80. A project costs 60. The safe policy produces estate 100;
the risky policy produces estate 20 or 220 with equal probabilities. Both pricing
and shareholder valuation use the stated risk-neutral probabilities and gross
discount factor $21/20$. Absolute priority in this illustration gives receipts
\begin{equation}
(s,j,e)=\bigl(\min(A,40),\min((A-40)^+,80),(A-120)^+\bigr).
\end{equation}
Thus the safe estate pays $(40,60,0)$, and the risky estates pay $(20,0,0)$
and $(40,80,100)$. Each allocation sums to its estate. The candidate junior
proceeds are $400/7$ under the safe policy and $800/21$ under the risky policy.

At a posted quote $P$, existing shareholders must subscribe $60-P$ to cover
the investment. Their net value from policy $\pi$ is
$\E[e_\pi]/(21/20)-(60-P)$. They prefer risk at either quoted price.
Consequently, the safe-policy quote is inconsistent with their best response:
if risk is chosen, junior lenders receive present value $800/21$ after advancing
$400/7$, a shortfall of $400/21$. The risky quote is consistent within the
pure-policy class, but requires shareholder cash $460/21$. If shareholders
have only 20, the project remains unfunded. The safe quote would need less cash,
but its incentive problem remains. Even an enforceable safe-policy commitment
would still need shareholder participation: here its zero terminal equity
receipt does not cover their positive subscription. Liquidity, incentives,
lender valuation and voluntary participation are distinct restrictions.

The same separation applies to restructuring. Given estate 70 and creditor
outside values 50 and 20, an offer 45/25 conserves value but the senior creditor
refuses. An offer 50/20 meets both constraints. A mandatory fee 1 reduces the
available estate to69 and makes simultaneous participation impossible under
these stated outside values. This is a feasibility result. Selecting a unique
offer requires bargaining rights or another decision rule; it does not follow
from total value conservation.

\section{Why positive surplus can still end in liquidation}
Hart and Moore's retained 1991 working paper makes the distinction between
productive value and credible repayment particularly clear. The entrepreneur
can withdraw human capital, the creditor can seize specified assets, and
renegotiation follows an explicit offer game. Returns and liquidation values
are deterministic, private savings have a particular seizure treatment, and
one creditor controls the assets \cite[Sections~2--3 and Appendix~A1--A4]{litclosehartmoore}.
Under those assumptions, Lemma~1 derives the creditor's continuation reservation
value and requires both future nonrepudiation and cash feasibility. The bargaining
rule determines a result that cannot be obtained merely by balancing accounts.

In their Example~1, take two remaining operating receipts of 2 and 26 and
liquidation rents of 8 in each period. The creditor's reservation value is
16, while the final credible payment is at most 13. The entrepreneur therefore
needs to provide at least 3 before the final date. The first operating receipt
provides only 2, so private wealth $w\ge1$ is necessary. At $w=1$, an immediate
payment of 1 followed by payments 2 and 13 satisfies the cash and repayment
conditions. At $w<1$, it cannot. Liquidation occurs despite operating value
28 exceeding liquidation value 16 \cite[Example~1, printed p.~16]{litclosehartmoore}.
The reference calculation reproduces this boundary with exact rational amounts.
It is a useful counterexample to assuming that positive total surplus guarantees
refinancing or restructuring.

Their repayment-path analysis also reinforces the distinction between an
observed payment and its contractual explanation. A long-term contract and
rolling short contracts can support the same path under the paper's conditions;
the latter may additionally require access to outside refinancing
\cite[Proposition~2 and printed pp.~27--28]{litclosehartmoore}. A future market closure
can therefore distinguish histories that looked equivalent before the shock.
For a multi-creditor company, the analyst must additionally specify priority,
control and consent rights, collateral and intercreditor restrictions. The
single-creditor result supplies neither an equity-issuance model nor a general
legal restructuring procedure.
